\section{The nine orders compared}
\label{app:orders}

The names below follow the usual Latin-derived English terminology. Rank~1 is the highest. Dionysius presents the triads in CH VI--IX; Gregory gives his ascending sequence in \emph{Homilies on the Gospels} XXXIV.7 and explains the offices in XXXIV.8--10. Aquinas compares the two arrangements in \ST{I q.~108 a.~6}. The biblical passages attest the names; their ordered synthesis belongs to the patristic exposition.

\begingroup
\small
\setlength{\tabcolsep}{4pt}
\begin{longtable}{@{}>{\raggedright\arraybackslash}p{.19\linewidth}>{\raggedright\arraybackslash}p{.16\linewidth}>{\raggedright\arraybackslash}p{.09\linewidth}>{\raggedright\arraybackslash}p{.09\linewidth}>{\raggedright\arraybackslash}p{.34\linewidth}@{}}
\toprule
Order & Latin & Dionysius & Gregory & Scriptural name \\
\midrule
\endhead
Seraphim & Seraphim & 1 & 1 & Isaiah 6:2, 6 \\
Cherubim & Cherubim & 2 & 2 & Genesis 3:24; Ezekiel 10 \\
Thrones & Throni & 3 & 3 & Colossians 1:16 \\
Dominions & Dominationes & 4 & 4 & Colossians 1:16; Ephesians 1:21 \\
Virtues & Virtutes & 5 & 7 & Ephesians 1:21 \\
Powers & Potestates & 6 & 6 & Colossians 1:16; Ephesians 1:21 \\
Principalities & Principatus & 7 & 5 & Colossians 1:16; Ephesians 1:21 \\
Archangels & Archangeli & 8 & 8 & Jude 9 \\
Angels & Angeli & 9 & 9 & Luke 1:26; Hebrews 1:14 \\
\bottomrule
\end{longtable}
\endgroup

In Dionysius the highest triad is seraphim, cherubim, and thrones; the middle is dominions, virtues, and powers; the last is principalities, archangels, and angels. Gregory exchanges the relative places of principalities and virtues. Aquinas explains how each arrangement follows the offices assigned to its terms. He preserves the names and their differences while showing a substantial agreement in the ordering of divine contemplation, governance, and execution.

\begingroup
\small
\setlength{\tabcolsep}{4pt}
\begin{longtable}{@{}>{\raggedright\arraybackslash}p{.18\linewidth}>{\raggedright\arraybackslash}p{.36\linewidth}>{\raggedright\arraybackslash}p{.39\linewidth}@{}}
\toprule
Order & Dionysian emphasis & Gregorian emphasis \\
\midrule
\endhead
Seraphim & Burning, purifying love; unceasing movement toward God (VII.1). & Incomparable ardor from nearness to the Creator (XXXIV.10). \\
Cherubim & Fullness of knowledge; contemplation and communication of wisdom (VII.1). & Fullness of knowledge through contemplation of divine glory (XXXIV.10). \\
Thrones & Stable receptivity and readiness to bear God (VII.1). & God exercises judgment through these seats of his presence (XXXIV.10). \\
Dominions & Freedom from servility; participation in divine lordship (VIII.1). & Eminent authority over subordinate ranks (XXXIV.10). \\
Virtues & Unwearied strength and readiness for divine action (VIII.1). & Ministers through whom signs and miracles are worked (XXXIV.10). \\
Powers & Ordered exercise of authority toward God (VIII.1). & Restraint of hostile spirits and their assaults (XXXIV.10). \\
Principalities & Godward leadership of subordinate ministers (IX.1). & Direction of good angels in the execution of divine ministries (XXXIV.10). \\
Archangels & Mediation between principalities and angels (IX.2). & Announcement of greater mysteries (XXXIV.8). \\
Angels & Completion of heavenly mediation toward humanity (IX.2). & Announcement of lesser matters; messenger as an office (XXXIV.8). \\
\bottomrule
\end{longtable}
\endgroup

Parker's wording must be translated consistently between the two tables and his chapter~VIII: \emph{Lordships} means dominions (\emph{kyriotetes}); \emph{Powers} means virtues (\emph{dynameis}); \emph{Authorities} means powers (\emph{exousiai}). The Latin \emph{virtus} here signifies strength or power. It does not designate a separate company composed of moral habits. Common names can also designate a perfection shared by all angels; CH V and XI distinguish that use from the proper name of an order.
